\section{La línea de supervivencia de una empresa emergente: guerra de subsidios, retiro de la plataforma y 4,000 yuanes en caja}

La sección anterior trató el producto; ahora pasamos al proceso de emprender, porque en una aplicación de IA el criterio de producto y la presión por sobrevivir van unidos. La descripción del video menciona la guerra de subsidios, el retiro del producto de la plataforma, que en la cuenta solo quedaban 4,000 yuanes en efectivo y que no había forma de conseguir financiamiento. La entrevista tiene una carga emocional fuerte porque el despegue de Lovart no fue un crecimiento gradual, sino un giro repentino que llegó después de una serie de reveses.

\begin{figure}[H]
\centering
\includegraphics[width=0.96\textwidth]{startup-survival-path.png}
\caption{La línea de supervivencia de una aplicación emergente: seguir con vida después de la guerra de subsidios, el retiro de la plataforma, el agotamiento del efectivo y la dificultad para conseguir financiamiento. Diagrama conceptual propio, elaborado a partir de la descripción del video y de la conversación entre 00:36:00 y 00:45:00.}
\end{figure}

\begin{knowledgebox}{Lectura de la figura: los momentos difíciles no son una anécdota, sino parte del criterio de producto}
La guerra de subsidios le mostró al equipo que el tráfico no es una barrera de entrada; el retiro de la plataforma le hizo ver el riesgo de depender de una plataforma; el agotamiento del efectivo lo obligó a concentrarse, y la imposibilidad de conseguir financiamiento puso a prueba la convicción del fundador. Todo esto influye en la ruta del producto.
\end{knowledgebox}

\begin{knowledgebox}{Tabla de riesgos de una empresa emergente de aplicaciones de IA}
\begin{tabularx}{\textwidth}{p{0.22\textwidth}p{0.34\textwidth}X}
\toprule
Riesgo & Manifestación & Respuesta \\
\midrule
Guerra de subsidios & Se compran usuarios con dinero y la retención no necesariamente es real & Encontrar cuanto antes la difusión orgánica y los casos de uso por los que se paga. \\
Retiro de la plataforma & Cambian las reglas de una plataforma externa o el producto se ve obligado a salir de ella & Conservar un acceso propio y la relación directa con los usuarios. \\
Efectivo agotado & Los fondos en la cuenta no alcanzan para seguir probando & Concentrarse en las funciones de mayor valor y extender el runway. \\
Dificultad para conseguir financiamiento & El mercado no cree en la ventana de oportunidad de las aplicaciones & Demostrar el PMF con usuarios e ingresos reales. \\
\bottomrule
\end{tabularx}
\end{knowledgebox}

\begin{knowledgebox}{Cómo influyen los momentos difíciles en la ruta del producto}
\begin{tabularx}{\textwidth}{p{0.22\textwidth}p{0.34\textwidth}X}
\toprule
Momento difícil & Lo que obliga al equipo a ver & Acción correcta que puede derivarse \\
\midrule
Guerra de subsidios & El crecimiento comprado no es confiable & Volver a la demanda orgánica y a la retención real. \\
Retiro de la plataforma & El acceso a través de una plataforma no se controla & Construir canales propios y respaldos en varias plataformas. \\
Efectivo agotado & El tiempo para probar es limitado & Concentrarse en el caso de uso más fuerte y en la ruta más corta para generar ingresos. \\
Sin financiamiento & El relato no convence a los inversionistas & Demostrarlo al revés, con usuarios e ingresos. \\
\bottomrule
\end{tabularx}
\end{knowledgebox}

\subsection{El CEO combativo}

Lo anterior trató la presión externa; esta sección examina cómo esa presión recae en el fundador. La descripción del video dice que Chen Mian (陈冕) se parece más a un CEO combativo. No se trata de una etiqueta de personalidad, sino de la realidad de emprender con aplicaciones de IA: el producto, el financiamiento, el crecimiento, la entrega y el estado de ánimo del equipo recaen en una sola persona. Las capacidades de los modelos cambian rápido y la ventana de oportunidad es corta, así que el fundador tiene que decidir rápido y bajo presión.

\begin{figure}[H]
\centering
\includegraphics[width=0.90\textwidth]{fighting-ceo.png}
\caption{El CEO combativo: producto, financiamiento, crecimiento, entrega y energía emocional, todo al máximo al mismo tiempo. Diagrama conceptual propio, elaborado a partir de la conversación entre 00:02:00 y 00:03:10 y entre 01:19:00 y 01:28:00.}
\end{figure}

\begin{warningbox}{Mucha energía emocional no es un sistema de gestión}
El entusiasmo y la combatividad del fundador pueden servir para atravesar los momentos difíciles, pero al final la empresa sigue necesitando procesos, organización y disciplina financiera. La alegría puede dar energía, pero no sustituye a la gestión.
\end{warningbox}

\subsection{Resumen de la sección}

La historia de Lovart muestra que la ventana de oportunidad de las aplicaciones de IA es muy corta y los momentos difíciles son muy duros. Sobrevivir hasta que la ventana se abre ya es, en sí, parte de la capacidad.
